\documentclass[12pt,letterpaper]{article}
\usepackage{graphicx,textcomp}
\usepackage{natbib}
\usepackage{setspace}
\usepackage{fullpage}
\usepackage{color}
\usepackage[reqno]{amsmath}
\usepackage{amsthm}
\usepackage{fancyvrb}
\usepackage{amssymb,enumerate}
\usepackage[all]{xy}
\usepackage{endnotes}
\usepackage{lscape}
\newtheorem{com}{Comment}
\usepackage{float}
\usepackage{hyperref}
\newtheorem{lem} {Lemma}
\newtheorem{prop}{Proposition}
\newtheorem{thm}{Theorem}
\newtheorem{defn}{Definition}
\newtheorem{cor}{Corollary}
\newtheorem{obs}{Observation}
\usepackage[compact]{titlesec}
\usepackage{dcolumn}
\usepackage{tikz}
\usetikzlibrary{arrows}
\usepackage{multirow}
\usepackage{xcolor}
\newcolumntype{.}{D{.}{.}{-1}}
\newcolumntype{d}[1]{D{.}{.}{#1}}
\definecolor{light-gray}{gray}{0.65}
\usepackage{url}
\usepackage{listings}
\usepackage{color}

\definecolor{codegreen}{rgb}{0,0.6,0}
\definecolor{codegray}{rgb}{0.5,0.5,0.5}
\definecolor{codepurple}{rgb}{0.58,0,0.82}
\definecolor{backcolour}{rgb}{0.95,0.95,0.92}

\lstdefinestyle{mystyle}{
	backgroundcolor=\color{backcolour},   
	commentstyle=\color{codegreen},
	keywordstyle=\color{magenta},
	numberstyle=\tiny\color{codegray},
	stringstyle=\color{codepurple},
	basicstyle=\footnotesize,
	breakatwhitespace=false,         
	breaklines=true,                 
	captionpos=b,                    
	keepspaces=true,                 
	numbers=left,                    
	numbersep=5pt,                  
	showspaces=false,                
	showstringspaces=false,
	showtabs=false,                  
	tabsize=2
}
\lstset{style=mystyle}
\newcommand{\Sref}[1]{Section~\ref{#1}}
\newtheorem{hyp}{Hypothesis}

\title{Problem Set 2}
\date{Due: February 18, 2026}
\author{Rosalie LaCroix}


\begin{document}
	\maketitle
	\section*{Instructions}
	\begin{itemize}
		\item Please show your work! You may lose points by simply writing in the answer. If the problem requires you to execute commands in \texttt{R}, please include the code you used to get your answers. Please also include the \texttt{.R} file that contains your code. If you are not sure if work needs to be shown for a particular problem, please ask.
		\item Your homework should be submitted electronically on GitHub in \texttt{.pdf} form.
		\item This problem set is due before 23:59 on Wednesday February 18, 2026. No late assignments will be accepted.

	\end{itemize}

	\vspace{.25cm}

\noindent We're interested in what types of international environmental agreements or policies people support (\href{https://www.pnas.org/content/110/34/13763}{Bechtel and Scheve 2013)}. So, we asked 8,500 individuals whether they support a given policy, and for each participant, we vary the (1) number of countries that participate in the international agreement and (2) sanctions for not following the agreement. \\

\noindent Load in the data labeled \texttt{climateSupport.RData} on GitHub, which contains an observational study of 8,500 observations.

\begin{itemize}
	\item
	Response variable: 
	\begin{itemize}
		\item \texttt{choice}: 1 if the individual agreed with the policy; 0 if the individual did not support the policy
	\end{itemize}
	\item
	Explanatory variables: 
	\begin{itemize}
		\item
		\texttt{countries}: Number of participating countries [20 of 192; 80 of 192; 160 of 192]
		\item
		\texttt{sanctions}: Sanctions for missing emission reduction targets [None, 5\%, 15\%, and 20\% of the monthly household costs given 2\% GDP growth]
		
	\end{itemize}
	
\end{itemize}

\newpage
\noindent Please answer the following questions:

\begin{enumerate}
	\item
	Remember, we are interested in predicting the likelihood of an individual supporting a policy based on the number of countries participating and the possible sanctions for non-compliance.
	\begin{enumerate}
		\item [] Fit an additive model. Provide the summary output, the global null hypothesis, and $p$-value. Please describe the results and provide a conclusion.

	\end{enumerate}
	
	\lstinputlisting[language=R, firstline=61, lastline=70]{PS02_RL.R}
	
\begin{table}[!htbp] \centering 
	\caption{Table of Coefficients} 
	\label{} 
	\begin{tabular}{@{\extracolsep{5pt}}lc} 
		\\[-1.8ex]\hline 
		\hline \\[-1.8ex] 
		& \multicolumn{1}{c}{\textit{Dependent variable:}} \\ 
		\cline{2-2} 
		\\[-1.8ex] & Choice \\ 
		& Additive Model \\ 
		\hline \\[-1.8ex] 
		Countries (80 of 192) & 0.336$^{***}$ \\ 
		& (0.054) \\ 
		& \\ 
		Countries (160 of 192) & 0.648$^{***}$ \\ 
		& (0.054) \\ 
		& \\ 
		Sanctions (5\%) & 0.192$^{***}$ \\ 
		& (0.062) \\ 
		& \\ 
		Sanctions (15\%) & $-$0.133$^{**}$ \\ 
		& (0.062) \\ 
		& \\ 
		Sanctions (20\%) & $-$0.304$^{***}$ \\ 
		& (0.062) \\ 
		& \\ 
		Constant & $-$0.273$^{***}$ \\ 
		& (0.054) \\ 
		& \\ 
		\hline \\[-1.8ex] 
		Observations & 8,500 \\ 
		Log Likelihood & $-$5,784.130 \\ 
		Akaike Inf. Crit. & 11,580.260 \\ 
		\hline 
		\hline \\[-1.8ex] 
		\textit{Note:}  & \multicolumn{1}{r}{$^{*}$p$<$0.1; $^{**}$p$<$0.05; $^{***}$p$<$0.01} \\ 
	\end{tabular} 
		\label{tab:Table1}
	\end{table} 
	
	\noindent The global null hypothesis states that neither of the variables, number of participating countries and sanctions for missing emission reduction targets, are significant in predicting whether or not an individual agrees with the given policy. The alternative hypothesis states that at least one of the variables listed is significant in predicting whether or not an individual agrees with the given policy. 
	
	$H_{0}$: $\beta_{1}$ = $\beta_{2}$ = 0, $H_{a}$ : At least one $\beta$ $\neq$ 0
	
	\noindent In \autoref{tab:Table1}, it can be seen that the coefficient for the number of participating countries, 0.458, is significant in predicting whether or not an individual agrees with a given policy. It can also be noted that the \autoref{tab:Table1} indicates that the coefficients for the impact of sanctions also have significant linear predictive power in determining whether or not an individual agrees with a certain policy. At this point in analysis, it can be determined that the null hypothesis that all of these coefficients are equal to zero can be rejected, and that all of them bar the quadratic relationship for the numeber of participating countries, are significant in predicting the support of a policy and should be included in the regression model.
	
\newpage
\item
If any of the explanatory variables are statistically significant in this model, then:
\begin{enumerate}
	\item
	For the policy in which nearly all countries participate [160 of 192], how does increasing sanctions from 5\% to 15\% change the odds that an individual will support the policy? (Interpretation of a coefficient)
	\vspace{.25cm}
		
	\noindent As an additive model has been fitted, the number of participating countries does not have an impact on the change in odds when increasing sanctions from 5\% to 15\%. The sanctions variable is in reference to the "None" level, so the difference between the coefficients for 5\% and 15\% sanctions must be mathematically compared against one another.
		
	\lstinputlisting[language=R, firstline=85, lastline=86]{PS02_RL.R}

	\begin{verbatim}
		 sanctions15% 
		 -0.3251028
	\end{verbatim}
			
	\noindent This indicates that the in accordance with our additive model, holding all other variables constant, when the sanctions are increased from 5\% to 15\%, the log odds that an individual will support the policy decrease by 0.325. 
	\vspace{.25cm}
		
	\item
	For the policy in which very few countries participate [20 of 192], how does increasing sanctions from 5\% to 15\% change the odds that an individual will support the policy? (Interpretation of a coefficient)
	\vspace{.25cm}
			
	\noindent As this is an additive model, the number of participating countries does not have an impact on the change in odds when increasing sanctions from 5\% to 15\%. To interpret the change in odds for increasing the sanctions, it will be the same at all levels of country participation and is therefore the same as in question 2(a). When sanctions are increased from 5\% to 15\%, and all other variables are held constant, the log odds that an individual will support the policy decrease by 0.325. 
	\vspace{.25cm}
			
	\item
	What is the estimated probability that an individual will support a policy if there are 80 of 192 countries participating with no sanctions? 
	\vspace{.25cm}
			
	\noindent The estimated probability that an individual will support a policy if there are 80 out of 192 countries participating with no sanctions can be determined by solving for the log odds and placing this into a formula to determine the probability that y = 1. The formula for this is as follows:
	
	\[\pi_{i} = P(Y_{i} = 1| X_{i}) = \frac{exp(\beta_{0} + \beta_{1}X_{1i} + ... \beta_{k}X_{ki})}{1 + exp(\beta_{0} + \beta_{1}X_{1i} + ... \beta_{k}X_{ki})}\]
	
	\noindent The coefficients for $\beta_{0}$ and $\beta_{1}$ can be input and utilized to solve for $\pi_{i}$, which is the probability that $Y_{i} = 1$ when $\beta_{1}$ represents the effect on the log odds that an individual will support a policy if there are 80 of 192 countries participating.
		
	\lstinputlisting[language=R, firstline=88, lastline=93]{PS02_RL.R}
	
	\begin{verbatim}
		0.5159191 
	\end{verbatim}
		
	\noindent The estimated probability that an individual will support a policy if there are 80 of 192 countries participating with no sanctions is 0.5159, or 51.59\%.
		
	\end{enumerate}
	\item
	Would the answers to 2a and 2b potentially change if we included an interaction term in this model? Why? 
	
	\noindent Yes, the answers to 2a and 2b would potentially change if we included an interaction term in this model. This is due to the fact that we would no longer have an additive model and would now have an interaction model. An interaction model would mean that there is an interaction effect between the explanatory variables, such that a change in one variable impacts how the other variable affects the log-odds of the outcome variable. In this specific data, an interaction term would mean that the number of countries participating impacts the effect of how the sanctions for missing reduction targets affects the log-odds an individual supports a policy, or vice versa. For this reason, the answers to 2a and 2b would potentially change, as the change in odds that an individual supports the policy due to increasing sanctions from 5\% to 15\% could be different at differing levels of country participation due to the interaction between country participation and sanctions.
	
	\begin{itemize}
		\item Perform a test to see if including an interaction is appropriate.
		
		\noindent 
		
		\lstinputlisting[language=R, firstline=100, lastline=103]{PS02_RL.R}
			
		\lstinputlisting[language=R, firstline=106, lastline=107]{PS02_RL.R}
		
	\begin{verbatim}
		Analysis of Deviance Table
		
		Model 1: choice ~ countries + sanctions
		Model 2: choice ~ countries * sanctions
		Resid. Df Resid. Dev Df Deviance Pr(>Chi)
		1      8494      11568                     
		2      8488      11562  6   6.2928   0.3912
	\end{verbatim}
		
	\end{itemize}
	\end{enumerate}


\end{document}
